%# -*- coding: utf-8-unix -*-
\chapter*{NEW DEVELOPMENTS}
\addcontentsline{toc}{chapter}{New Developments}
\markboth{NEW DEVELOPMENTS}{NEW DEVELOPMENTS}

We give a brief summary of some of the progress in transcendental number\index{transcendental number} theory that has been made since this book was written. It has in fact been a very active field of research and only a few of the new results will be mentioned here. A full account of much of the ground that has been won recently can be found in \textit{Transcendence theory: advances and applications} (Academic Press, London and New York, 1977), Proceedings of a conference held in Cambridge in 1976, edited by A. Baker\index{Baker, R. C.}\index{Baker, A.} and D. W. Masser\index{Masser, D. W.}; it will be referred to briefly as TTAA.

\section*{1. Linear forms in logarithms\index{linear forms in logarithms}}

The inequalities for \(\Lambda\) recorded after the enunciation of \autoref{thm:ch3_1} have been much improved. It is shown in Chapter 1 of TTAA that if \(\alpha_i\) and \(\beta_i\) are algebraic numbers\index{algebraic number} with heights at most \(A_i\) (\(\ge 4\)) and \(B\) (\(\ge 4\)), and if the field generated by the \(\alpha\)'s and \(\beta\)'s over the rationals has degree at most \(d\), then either \(\Lambda = 0\) or \(|\Lambda| > (B\Omega)^{-C\Omega \log \Omega'}\), where \(\Omega = \log A_1 \dots \log A_n\), \(\Omega' = \Omega/\log A_n\) and \(C = (16nd)^{200n}\). Furthermore, in the rational case, that is when \(\beta_0 = 0\) and \(\beta_1, \dots, \beta_n\) are rational integers, the bracketed factor \(\Omega\) has been eliminated to yield \(|\Lambda| > B^{-C\Omega \log \Omega'}\). The latter result is best possible with respect to \(A_n\) and \(B\), and likewise with respect to \(A_1, \dots, A_{n-1}\) except for the second-order term \(\log \Omega'\). The proof involves, in particular, an improved version, due to Tijdeman\index{Tijdeman, R.}, of \autoref{lem:ch2_1} on page 25, a new idea of Shorey\index{Shorey, T. N.} on the size of the inductive steps which occur in \autoref{lem:ch2_6} on page 31, and some rather deep developments concerning the argument on pages 34 and 35 involving more widespread use of Kummer theory\index{Kummer theory}. Refinements in the results, relating especially to the expression for \(C\), have been announced by several writers.\footnote{cf. J. Australian Math. Soc. 25 (1978), 445-78.}

There is also an extensive \(p\)-adic theory of linear forms in logarithms\index{linear forms in logarithms}. The subject was initiated by Mahler\index{Mahler, K.} in the 1930s when he obtained \(p\)-adic analogues of both the Hermite\index{Hermite, Ch.}-Lindemann\index{Lindemann, F.} and the Gelfond\index{Gelfond, A. O.}-Schneider\index{Schneider, Th.} theorems; in fact, in the course of the work, Mahler\index{Mahler, K.} laid the foundations of the \(p\)-adic theory of analytic functions that has been fundamental to all later studies in the field. A good account of the subject is given by van der Poorten\index{van der Poorten, A. J.} in Chapter 2 of TTAA; in particular he establishes there an estimate for the \(p\)-adic valuation of \(\Lambda\) of essentially the same degree of precision as that quoted earlier in the Archimedean case.

\section*{2. Diophantine equations\index{Diophantine equations}}

The sharpened estimates for \(\Lambda\) and, in particular, their best possible dependence on \(A_n\), have led to some remarkable developments in connection with the theory of Diophantine equations\index{Diophantine equations}.

First, they yield at once an explicit upper bound for all solutions in integers \(x>1\), \(y>1\), \(n>2\) of the equation \(ax^n-by^n=c\), where \(a,b,c\) (\(\neq 0\)) are any given integers. In fact, assuming, as we may without loss of generality, that \(a\) and \(b\) are positive, and that \(y\ge x\), the equation gives \(|\Lambda|\ll y^{-n}\), where \(\Lambda=\log{(a/b)}+n\log{(x/y)}\), and the implied constant depends only on \(a\), \(b\) and \(c\). But, from the results recorded in \S~1, we have \(\log{|\Lambda|} \gg -\log{y}\log{n}\); a comparison of estimates yields a bound for \(n\) in terms of \(a\), \(b\), \(c\), and \autoref{thm:ch4_2} then furnishes bounds for \(x\) and \(y\).

A similar argument was employed by Tijdeman\index{Tijdeman, R.}\footnote{Acta Arith. 29 (1976), 197-209.} to show that the famous Catalan equation\index{Catalan equation} \(x^p - y^q = 1\) has only finitely many solutions in integers \(x, y, p, q\) (all \({}> 1\)). Here we can assume that \(p, q\) are odd primes, and it will suffice to treat the case \(p > q\). Then \(x = kX^q + 1\), \(y = lY^p - 1\) for some integers \(X, Y\), where \(k\) is \(1\) or \(1/p\) and \(l\) is \(1\) or \(1/q\). Plainly we have
\[ |p\log x - q\log y| \ll y^{-q}, \]
and thus \(|\Lambda| \ll (kX^q)^{-1}\), where
\[ \Lambda = p \log k - q \log l + pq \log (X/Y). \]
Hence, by \S~1, we obtain \(q \ll (\log p)^4\). Further, the linear form \(\Lambda' = p \log (x/Y^q) - q \log l\) satisfies \(|\Lambda'| \ll (lY^p)^{-1}\), and so again by \S~1, we have \(p \ll q (\log p)^3\). Thus \(p\) and \(q\) are bounded, and so, by \autoref{thm:ch4_2}, also \(x\) and \(y\) are bounded, as required.

Many novel results have been obtained subsequent to Tijdeman\index{Tijdeman, R.}'s discovery; see especially Chapter 3 of TTAA. In particular, Schinzel\index{Schinzel, A.} and Tijdeman\index{Tijdeman, R.}\footnote{Acta Arith. 31 (1976), 199-204.} have effectively resolved the equation \(y^m = f(x)\), of hyperelliptic type, in integers \(x, y > 1, m > 2\). Shorey\index{Shorey, T. N.} and Tijdeman\index{Tijdeman, R.}\footnote{Math. Scand. 39 (1976), 5-18.} have dealt with the equation \(y^m = x^n + x^{n-1} + \dots + 1\) in integers \(x, y, m, n\) (all \({}>1\)) when either \(x\) is fixed or \(n + 1\) or \(y\) has a fixed prime factor. Gy{\"o}ry, Tijdeman\index{Tijdeman, R.} and Voorhoeve\index{Voorhoeve, M.}\footnote{Acta Arith. 38 (1981), 389-410.} have shown that for any integer \(k \ge 6\) the equation \(y^m = 1^k + \dots + x^k\) has only finitely many solutions in integers \(x \ge 1\), \(y \ge 1\), \(m > 1\) all of which can, in principle, be effectively determined. Stewart\index{Stewart, C. L.}\footnote{Mathematika 24 (1977), 130-2. A similar result was obtained independently by Inkeri and van der Poorten\index{van der Poorten, A. J.}.} has proved that the Fermat equation\index{Fermat equation} \(x^n + y^n = z^n\) has only finitely many solutions in positive integers \(x, y, z\) and \(n\) (\({}>2\)), provided that \(|x-y|\) is bounded. Van der Poorten\footnote{Acta Arith. 33 (1977), 195-207.} has used \(p\)-adic analysis to effectively resolve the equation \(x^p - y^q = z^r\) in integers \(x, y, z, p, q\) (all \({}> 1\)), where \(z\) is composed solely of powers of a fixed set of primes, and \(r\) is the least common multiple of \(p\) and \(q\). And Gy{\"o}ry and Papp\index{Papp, Z. Z.}\footnote{Publ. Math. Debrecen 25 (1978), 311-25.} have recently made valuable progress in connection with the problem of effectively solving equations of norm form in several variables of the kind referred to on page 68 of \autoref{ch:7}.

It should also be noted that the work of Schmidt\index{Schmidt, W. M.} discussed in \autoref{ch:7} has been extensively generalized in the \(p\)-adic domain by Dubois\index{Dubois, E.} and Rhin\index{Rhin, G.},\footnote{Ast{\'e}risque (Soc. Math. France), 24-5 (1975), 211-27.} and by Schlickewei\index{Schlickewei, H. P.}.\footnote{J.M. 288 (1976), 86-105.}

\section*{3. Elliptic and Abelian functions\index{Abelian function}}

Much research has been carried out in the last few years on the work of \autoref{ch:6}. In fact the results have been generalized in three directions, namely quantitatively, \(p\)-adically and in the context of Abelian varieties. Of particular interest are quantitative versions of Masser\index{Masser, D. W.}'s result cited at the end of page 64; see the paper by Anderson\index{Anderson, M.} in Chapter 7 of TTAA. The studies lead to delicate questions on the division value properties of the Weierstrass\index{Weierstrass, K.} functions, and, in particular, some recent theorems of Bashmakov\index{Bashmakov, M.} and Ribet\index{Ribet, K. A.} on the subject have proved useful. The methods can be generalized, as Masser\index{Masser, D. W.},\footnote{Invent. Math. 45 (1978), 61--82.} Coates\index{Coates, J.} and Lang\index{Lang, S.}\footnote{Invent. Math. 34 (1976), 129--133.} have shown, to deal with Abelian functions\index{Abelian function}, provided that one makes suitable assumptions concerning complex multiplications\index{complex multiplication}, and the theory has been carried over to the \(p\)-adic domain for a wide class of primes \(p\) by Bertrand\index{Bertrand, D.} and Flicker\index{Flicker, Y. Z.}.\footnote{Acta Arith. 38 (1980), 47--61.}

Other noteworthy results are a \(p\)-adic version of Schneider\index{Schneider, Th.}'s \autoref{thm:ch6_2} obtained by Bertrand\index{Bertrand, D.},\footnote{Chapter 9 of TTAA, and Invent. Math. 40 (1977), 171-93.} a solution to the open problem mentioned on page 57 in the case \(p=2\) by Masser\index{Masser, D. W.},\footnote{Mathematika, 22 (1975), 97-107.} his extension of \autoref{thm:ch6_6} to include also the number \(2\pi i\),\footnote{Chapter 6 of TTAA.} and a striking theorem of Chudnovsky\index{Chudnovsky, G. V.}\footnote{See the article by Waldschmidt\index{Waldschmidt, M.} in Lecture Notes in Mathematics, vol. 567 (Springer-Verlag, Berlin, 1977), pp. 274-92.} to the effect that if \(g_2\) and \(g_3\) are algebraic then the transcendence degree\index{transcendence degree} of the field generated by \(\omega_1\), \(\omega_2\), \(\eta_1\), \(\eta_2\) over the rationals is at least 2. As an immediate corollary one obtains the transcendence of \(\Gamma(\frac{1}{4})\); this follows on considering the curve \(y^2 = 4x^3 - 4x\) for which one can take \(\omega_1 = (\Gamma(\frac{1}{4}))^2/\sqrt{(8\pi)}\), \(\omega_2 = i\omega_1\), \(\eta_1 = \pi/\omega_1\) and \(\eta_2 = -i\eta_1\). Similarly from the curve \(y^2 = 4x^3 - 4\) one deduces that \(\Gamma(\frac{1}{3})\) is transcendental.

\section*{4. Arithmetical properties of meromorphic functions\index{meromorphic function}}

An excellent account of the Siegel\index{Siegel, C. L.}-Shidlovsky\index{Shidlovsky, A. B.} theorems is given by Mahler\index{Mahler, K.} in his tract \textit{Lectures on transcendental numbers\index{transcendental number}\index{numbers}} (Lecture Notes in Mathematics, vol. 546, Springer-Verlag, Berlin, 1976). The volume contains also some notable examples of entire transcendental functions which, together with their derivatives, assume algebraic values at all algebraic points, and also an interesting appendix giving an historical survey of the classical proofs of the transcendence of \(e\) and \(\pi\). The papers of Waldschmidt\index{Waldschmidt, M.} and Bertrand\index{Bertrand, D.} in Chapters 11 and 12 of TTAA contain new results on meromorphic functions\index{meromorphic function} of finite order.

The problem of estimating the quantity \(P(E_1, \dots, E_n)\) referred to on page 117 has been studied by Nesterenko\index{Nesterenko, Ju. V.},\footnote{I.A.N. 41 (1977), 253-84.} and he has obtained bounds with an explicit dependence on the degree of \(P\); further work in this connection, and also in relation to \autoref{ch:10}, has been carried out by V{\"a}{\"a}n{\"a}nen.\footnote{Manuscripta Math. 21 (1977), 173-80; Bull. Australian Math. Soc. 14 (1976), 161-79.} The question of estimating similar expressions for \(G\)-functions has been investigated by Nurmagomedov\index{Nurmagomedov, M. S.}, Galo\v{c}kin and, in the \(p\)-adic domain, by Flicker\index{Flicker, Y. Z.}.\footnote{J. London Math. Soc. 15 (1977), 395-402.} Incidentally, no satisfactory \(p\)-adic generalization of the Siegel\index{Siegel, C. L.}-Shidlovsky\index{Shidlovsky, A. B.} theory has been given as yet.

In another direction, an old and previously neglected method of Mahler\index{Mahler, K.} on functions satisfying certain functional equations has become the focus of much activity in the last few years. The principles underlying the method are simpler than those of the Siegel\index{Siegel, C. L.}-Shidlovsky\index{Shidlovsky, A. B.} theory, and they have proved to be capable of considerable generalization; see the survey articles of Loxton\index{Loxton, J. H.} and van der Poorten\index{van der Poorten, A. J.} and of Kubota\index{Kubota, K. K.} in Chapters 15 and 16 of TTAA. Typical examples of the results furnished by the method are the transcendence of the Fredholm series\index{Fredholm series} \(\sum z^{2^n}\) for all algebraic \(z\) with \(0 < |z| < 1\), and of \(\sum [n\omega]z^n\) for all such \(z\) and any irrational \(\omega\). It also yields several new instances of algebraically independent numbers\index{numbers}.

\section*{5. Further works}

Stewart\index{Stewart, C. L.}\footnote{Acta Arith. 26 (1975), 427-33; Proc. London Math. Soc. 35 (1977), 425-47; see also Chapter 4 of TTAA.} has utilized the estimates for linear forms in logarithms\index{linear forms in logarithms} quoted earlier to give much new information on divisor properties of arithmetical sequences, and, in particular, on Lucas and Lehmer numbers\index{Lucas and Lehmer numbers}.

Chudnovsky\index{Chudnovsky, G. V.} has announced some extensive developments of the method of Gelfond\index{Gelfond, A. O.} discussed in \autoref{ch:12} leading to the algebraic independence of three or more numbers\index{numbers} from certain sets, rather than just two as occur in \autoref{thm:ch12_1} and \autoref{thm:ch12_2}.\footnote{See again the article by Waldschmidt\index{Waldschmidt, M.} in Lecture Notes in Mathematics, vol. 567 (Springer-Verlag, Berlin, 1977), pp. 274-92.}

R. C. Baker\index{Baker, R. C.}\index{Baker, A.}\footnote{Mathematika, 23 (1976), 18-31, 184-97.} has obtained new results on \(T\)-numbers, and also on a conjecture of Schmidt\index{Schmidt, W. M.} and the author concerning the metrical theory\index{metrical theory} discussed in \autoref{ch:10}.

Stepanov\index{Stepanov, S. A.}, and subsequently Bombieri\index{Bombieri, E.} and Schmidt\index{Schmidt, W. M.}, used transcendence techniques to furnish a new proof of the Riemann hypothesis\index{Riemann hypothesis} for curves which is of a much more elementary nature than that given originally.\footnote{See Schmidt\index{Schmidt, W. M.}'s tract, vol. 536 (1976) of Lecture Notes in Mathematics.}

And finally we mention that Ap{\'e}ry has just established the irrationality of \(\zeta(3)\).

\section*{6. Additional remarks (1990)}

A natural sequel to TTAA (see page 155) is \textit{New advances in transcendence theory} (Cambridge University Press, 1988), Proceedings of a Symposium held in Durham in 1986, edited by A. Baker\index{Baker, R. C.}\index{Baker, A.}; we shall refer to the volume briefly as NATT. It covers the most significant discoveries of recent years.

The papers in NATT by W{\"u}stholz and by Philippon and Waldschmidt\index{Waldschmidt, M.} establish, independently, that the second-order term \(\log \Omega'\) can be eliminated from the estimates for \(|\Lambda|\) discussed on page 155. The arguments depend on the theory of multiplicity estimates on group varieties. This is an important new instrument in transcendence originating from techniques involving commutative algebra introduced by Nesterenko\index{Nesterenko, Ju. V.};\footnote{I.A.N. 41 (1977), 253--284.} it has been developed by Brownawell\index{Brownawell, W.} and Masser,\footnote{J.M. 314 (1980), 200--216; Duke Math. J. 47 (1980), 273--295.} Masser\index{Masser, D. W.} and W{\"u}stholz,\footnote{Invent. Math. 64 (1981), 489--516; 80 (1985), 233--267.} Philippon\footnote{Bull. Soc. Math. France 114 (1986), 355--383.} and, most notably, by W{\"u}stholz.\footnote{J.M. 317 (1980), 102--119; Ann. Math. 129 (1989), 471--517.} The work has led to a remarkable synthesis of many aspects of the subject in terms of algebraic groups.

There are several substantial memoirs in NATT relating to Diophantine analysis\index{Diophantine analysis}; they include, in particular, an article by Evertse, Gy{\"o}ry, Stewart\index{Stewart, C. L.} and Tijdeman\index{Tijdeman, R.} surveying the fertile field of \(S\)-unit equations and a paper by Odoni furnishing an application to a question of Serre\index{Serre, J. P.} on modular forms. It is apparent, vide the tract\footnote{Here, much rests on the \(p\)-adic estimate of van der Poorten\index{van der Poorten, A. J.} mentioned on page 156. Yu Kunrui has recently given an improved version (Compositio Math., to appear) overcoming some errors in van der Poorten\index{van der Poorten, A. J.}'s demonstration; see the discussion in NATT.} \textit{Exponential Diophantine equations\index{Diophantine equations}}\footnote{Cambridge University Press, 1986; by T. N. Shorey\index{Shorey, T. N.} and R. Tijdeman\index{Tijdeman, R.}.} that the topic has become an unusually rich area of research during the past decade. In another exciting development, two new proofs of the Mordell\index{Mordell, L. J.} conjecture on the finiteness of the number of rational points on algebraic curves of genus\index{curves of genus} \({}>1\) (first established by Faltings in 1983) have recently been obtained by Masser\index{Masser, D. W.} and W{\"u}stholz and by Vojta using transcendence methods.

The study of \(E\)-functions and \(G\)-functions has continued to be a source of much activity. Again, a good account is provided by NATT; there are, especially, some beautiful results of Beukers\index{Beukers, F.} and Wolfart specifying precisely when the classical hypergeometric functions\index{hypergeometric function} can take algebraic values. A comprehensive exposition of the basic Siegel\index{Siegel, C. L.}-Shidlovsky\index{Shidlovsky, A. B.} theory is given in \textit{Transcendental numbers\index{transcendental number}\index{numbers}} (de Gruyter, 1989), Studies in mathematics 12, by A. B. Shidlovsky\index{Shidlovsky, A. B.}.

The metrical conjecture concerning \((\psi(h))^n\) mentioned on page 95 has been proved by V. I. Bernik; references to this and also to his solution to a problem of Baker\index{Baker, R. C.}\index{Baker, A.} and Schmidt\index{Schmidt, W. M.} on Hausdorff dimension\index{Hausdorff dimension} are given in his paper in NATT.

The famous problem of Gauss\index{Gauss, C. F.} of determining effectively all the imaginary quadratic fields with a given class number, earlier resolved for class numbers\index{numbers} 1 and 2 as described in \autoref{ch:5}, has now been resolved in general through the theory of elliptic curves\index{elliptic curve}, by Goldfeld, Gross and Zagier.\footnote{C.R. 297 (1983), 85--87; Bull. Amer. Math. Soc. 13 (1985), 23--37.}
